\section*{Capítulo \ref{ch:b1:earthquake-source} --- A fonte do terremoto}

\begin{solution}{exo:b1:earthquake-source:1}
(a) $\log_{10}(\num{4e17}) = 17.60$, logo $M_w = \tfrac23(17.60 - 9.1) = 5.67$.
(b) $\log_{10}(\num{5.59e22}) = 22.75$ e $M_w = \tfrac23(22.75 - 9.1) =
9.10$. (c) $M_0 = 10^{1.5\times7.0 + 9.1} = 10^{19.6} = \qty{3.98e19}{N.m}$.
\end{solution}

\begin{solution}{exo:b1:earthquake-source:2}
(a) $T = D/\dot s = \qty{3}{m}/\qty{0.025}{m.yr^{-1}} = \qty{120}{yr}$.
(b) Com $c = 0.5$, $T = D/(c\dot s) = \qty{240}{yr}$: o rastejamento absorve metade
do movimento, e os terremotos, do mesmo deslizamento, vêm com metade da frequência.
\end{solution}

\begin{solution}{exo:b1:earthquake-source:3}
Pela \cref{prop:b1:earthquake-source:moment-steps}, a razão dos momentos é
$10^{1.5\times(7.4 - 6.2)} = 10^{1.8} = 63.1$: são necessários 63 terremotos de magnitude 6.2 para liberar o momento de um de magnitude 7.4.
\end{solution}

\begin{solution}{exo:b1:earthquake-source:4}
Com $b = 1$, cada unidade de magnitude divide a contagem por 10: 10 terremotos de magnitude 5 ou mais por ano, 1 de magnitude 6 ou mais, e 0.1 de magnitude 7 ou mais, isto é, um a cada 10 anos em média.
\end{solution}

\begin{solution}{exo:b1:earthquake-source:5}
(a) $A = \num{4.3e5} \times \qty{1.5e4}{m} = \qty{6.45e9}{m^2}$, e
$M_0 = \num{3e10}\times\num{6.45e9}\times4 = \qty{7.74e20}{N.m}$;
$\log_{10}M_0 = 20.89$ e $M_w = \tfrac23(20.89 - 9.1) = 7.86$. (b) Uma profundidade de
\qty{10}{km} divide a área e o momento por 1.5, o que reduz a magnitude em
$\tfrac23\log_{10}1.5 = 0.12$, para 7.74. (c) A primeira, 7.86, arredonda para o 7.9 do catálogo; a segunda é um pouco baixa. A magnitude não é muito sensível à profundidade: um fator de 1.5 na área vale apenas 0.12 de magnitude.
\end{solution}

\begin{solution}{exo:b1:earthquake-source:6}
(a) $D = M_0/(\mu A) = \num{5.59e22}/(\num{3e10}\times\num{4e5}\times\num{1.5e5})
= \qty{31.1}{m}$. (b) $\qty{31.1}{m}/\qty{0.083}{m.yr^{-1}} = \qty{374}{yr}$.
(c) Com a fração travada 0.8, $31.1/(0.8\times0.083) = \qty{468}{yr}$; com
0.5, \qty{749}{yr}. Quanto menos travada a interface, mais tempo leva para se acumular o deslizamento de um grande terremoto; a média de deslizamento encontrada aqui com um retângulo de \qtyproduct{400 x 150}{km} e $\mu = \qty{30}{GPa}$ é ela mesma apenas uma estimativa, e o deslizamento real esteve longe de ser uniforme.
\end{solution}

\begin{solution}{exo:b1:earthquake-source:7}
(a) A \omterm{def:b1:geological-maps:attitude}{direção} é $\phi = 329$, o \omterm{def:b1:geological-maps:attitude}{sentido de mergulho} $a = 059$, o \omterm{def:b1:geological-maps:attitude}{mergulho} $\delta =
\ang{8}$. Então, em coordenadas, $\mathbf s = (\sin329, \cos329, 0) = (-0.5150, 0.8572, 0)$,
$\mathbf u = (-\cos8\sin59, -\cos8\cos59, \sin8) = (-0.8488, -0.5100, 0.1392)$,
$\mathbf n = (\sin8\sin59, \sin8\cos59, \cos8) = (0.1193, 0.0717, 0.9903)$,
e com $r = \ang{110}$,
\[
  \mathbf d = \cos110\,\mathbf s + \sin110\,\mathbf u = (-0.6215, -0.7724, 0.1308) .
\]
(b) O \omterm{def:b1:earthquake-source:focal-mechanism}{plano auxiliar} tem a normal $\mathbf d$ e o deslizamento $\mathbf n$.
Tomando a normal para cima, ele mergulha $\arccos 0.1308 = \ang{82.5}$ em direção a
$\arctan(-0.6215/-0.7724) = 218.8$, logo sua \omterm{def:b1:geological-maps:attitude}{direção} é $128.8$; o deslizamento
$\mathbf n$ projetado sobre seus vetores de \omterm{def:b1:geological-maps:attitude}{direção} e de maior declive dá o \omterm{def:b1:faults-folds-fabrics:fault-elements}{rake}
\ang{87.2}. Arredondando: $129/83/87$, o segundo plano da solução publicada.
(c) $\mathbf P = (\mathbf n - \mathbf d)/\sqrt2$ tem \omterm{def:b1:field-lab-methods-1:trend-plunge}{rumo} $221$ e \omterm{def:b1:field-lab-methods-1:trend-plunge}{caimento}
\ang{37}; $\mathbf T = (\mathbf n + \mathbf d)/\sqrt2$ tem \omterm{def:b1:field-lab-methods-1:trend-plunge}{rumo} $036$ e \omterm{def:b1:field-lab-methods-1:trend-plunge}{caimento}
\ang{52} (publicado: $222/38$ e $36/52$). (d) O plano raso, mergulhando
\ang{8} em direção a $059$, ao longo da laje subductante: uma \omterm{def:b1:faults-folds-fabrics:fault-types}{falha inversa} com \omterm{def:b1:faults-folds-fabrics:fault-elements}{rake} de
\ang{110} e, mergulhando menos que \ang{45}, um empurrão, a interface de placas da zona de subducção.
\end{solution}

\begin{solution}{exo:b1:earthquake-source:8}
(a) O plano mergulha em direção a $135$, logo $\mathbf n = (\sin135, \cos135, 0) =
(0.707, -0.707, 0)$; um deslizamento \omterm{def:b1:faults-folds-fabrics:fault-types}{dextral} tem $r = \ang{180}$, $\mathbf d = -\mathbf s
= (-0.707, -0.707, 0)$. Em \omterm{def:b1:geological-maps:attitude}{direção} ao norte, $\mathbf x = (0, 1, 0)$:
$(\mathbf n\cdot\mathbf x)(\mathbf d\cdot\mathbf x) = (-0.707)(-0.707) > 0$, um
empurrão. Leste: $(0.707)(-0.707) < 0$, um puxão. Sul: empurrão; oeste: puxão.
(b) $\mathbf P = (\mathbf n - \mathbf d)/\sqrt2 = (1, 0, 0)$: o \omterm{def:b1:earthquake-source:focal-mechanism}{eixo P} é
horizontal e tem \omterm{def:b1:field-lab-methods-1:trend-plunge}{rumo} $090$--$270$; $\mathbf T = (0, -1, 0)$ tem \omterm{def:b1:field-lab-methods-1:trend-plunge}{rumo}
$000$--$180$. (c) O plano vertical de \omterm{def:b1:geological-maps:attitude}{direção} $135$--$315$ tem a normal
$\mathbf d$.
Com o deslizamento $\mathbf n$, o bloco de seu lado sudoeste, para o qual $\mathbf d$ aponta, move-se em direção a $135$; visto do bloco nordeste, olhando para sudoeste, isso está à esquerda: um deslizamento \omterm{def:b1:faults-folds-fabrics:fault-types}{sinistral}. Pela \cref{prop:b1:earthquake-source:auxiliary}, este deslizamento dá os mesmos \omterm{def:b1:earthquake-source:focal-mechanism}{primeiros movimentos} em toda \omterm{def:b1:geological-maps:attitude}{direção}.
\end{solution}

\begin{solution}{exo:b1:earthquake-source:9}
(a) $V = \qty{40}{mm.yr^{-1}}$ e $v = V/4$ a \qty{10}{km}, logo pela
\cref{prop:b1:earthquake-source:profile}(b) $D = \qty{10}{km}$; diretamente,
$D = x/\tan(\pi v/V) = 10/\tan(\pi/4) = \qty{10}{km}$. (b) $v(30) =
(40/\pi)\arctan 3 = \qty{15.9}{mm.yr^{-1}}$. (c) $V/(\pi D) =
\num{0.040}/(\pi\times\num{10000}) = \num{1.27e-6}$ por ano.
\end{solution}

\begin{solution}{exo:b1:earthquake-source:10}
(a) Os momentos das faixas continuam $C q^k$ com $q = 10^{1.5 - b}$, mas agora $q < 1$: o momento das faixas decresce em direção às grandes magnitudes, e a faixa inferior, $k = 0$, leva a fração $C/S = (1 - q)/(1 - q^{K+1}) > 1 - q$, imagem espelhada de \cref{prop:b1:earthquake-source:largest}. Os menores terremotos contados dominariam então, e o total cresceria sem limite à medida que se incluíssem cada vez menores. (b) $1 - 10^{-0.7} = 0.80$: pelo menos \qty{80}{\percent}. (c) $\log_{10}N = 8.20 - 1.01\times9 = -0.89$, logo $N = 0.13$
por ano e cerca de 6.4 em cinquenta anos, três vezes os dois observados. A reta, ajustada entre as magnitudes 6 e 7.5, superestima os maiores terremotos; ou a lei se curva para baixo no topo, como sugere o tamanho finito das falhas, ou cinquenta anos são poucos para amostrar eventos que ocorrem uma vez a cada várias décadas. As duas possibilidades são compatíveis com dois eventos.
\end{solution}

\begin{solution}{exo:b1:earthquake-source:11}
(a) Em 150 anos um poste em $x$ se move $150\,v(x) = 150\times(34/\pi)
\arctan(x/15)$ milímetros em relação ao poste sobre a falha: \qty{0.11}{m} a \qty{1}{km}, \qty{1.28}{m} a \qty{15}{km}, \qty{2.08}{m} a \qty{50}{km}, em direção ao sul, para uma falha \omterm{def:b1:faults-folds-fabrics:fault-types}{dextral} vista do oeste. (b) O déficit de deslizamento é
$s = \qty{0.034}{m.yr^{-1}}\times\qty{150}{yr} = \qty{5.1}{m}$. Durante o terremoto os postes se movem $(s/\pi)\arctan(D/x)$: \qty{2.44}{m} a
\qty{1}{km}, \qty{1.28}{m} a \qty{15}{km}, \qty{0.47}{m} a \qty{50}{km}. (c) As somas são $0.11 + 2.44$, $1.28 + 1.28$ e $2.08 + 0.47$, todas
\qty{2.55}{m} $= s/2$: depois do terremoto a cerca é reta de novo do lado leste, deslocada de $s/2$, e o lado oeste, por simetria, de $-s/2$.
\end{solution}

\begin{solution}{exo:b1:earthquake-source:12}
(a) $\mathbf P = (0, 1, 0)$ e $\mathbf T = (1, 0, 0)$. O par
$\mathbf n = (\mathbf T + \mathbf P)/\sqrt2 = (0.707, 0.707, 0)$,
$\mathbf d = (\mathbf T - \mathbf P)/\sqrt2 = (0.707, -0.707, 0)$ satisfaz
$\mathbf P = (\mathbf n - \mathbf d)/\sqrt2$ e $\mathbf T = (\mathbf n +
\mathbf d)/\sqrt2$. Ambos os vetores são horizontais, logo ambos os \omterm{def:b1:earthquake-source:focal-mechanism}{planos nodais} são verticais: o plano de normal $\mathbf n$ (em direção a $045$) tem \omterm{def:b1:geological-maps:attitude}{direção} $135$--$315$, o plano de normal $\mathbf d$ (em direção a $135$) tem \omterm{def:b1:geological-maps:attitude}{direção} $045$--$225$. (b) No plano de \omterm{def:b1:geological-maps:attitude}{direção} $315$, o bloco do lado nordeste move-se ao longo de $\mathbf d$,
em direção a $135$: visto do outro bloco, à direita, um deslizamento \omterm{def:b1:faults-folds-fabrics:fault-types}{dextral}. No plano de \omterm{def:b1:geological-maps:attitude}{direção} $045$, o bloco do lado sudeste move-se ao longo de $\mathbf n$,
em direção a $045$: um deslizamento \omterm{def:b1:faults-folds-fabrics:fault-types}{sinistral}. (c) Uma falha nova faz $\ang{45} - \phi/2 =
\ang{30}$ com $\sigma_1$: com $\sigma_1$ de \omterm{def:b1:field-lab-methods-1:trend-plunge}{rumo} $000$, as falhas teriam \omterm{def:b1:geological-maps:attitude}{direção} $030$ (\omterm{def:b1:faults-folds-fabrics:fault-types}{sinistral}) e $330$ (\omterm{def:b1:faults-folds-fabrics:fault-types}{dextral}), a \ang{15} dos \omterm{def:b1:earthquake-source:focal-mechanism}{planos nodais}.
O \omterm{def:b1:earthquake-source:focal-mechanism}{eixo P} fica fixo a \ang{45} dos \omterm{def:b1:earthquake-source:focal-mechanism}{planos nodais} pela geometria, qualquer que seja o atrito da rocha, de modo que coincide com $\sigma_1$ apenas se a falha tiver deslizado a \ang{45} dele. (d) A falha de \omterm{def:b1:geological-maps:attitude}{direção} $040$ está a \ang{5} do \omterm{def:b1:earthquake-source:focal-mechanism}{plano nodal} de \omterm{def:b1:geological-maps:attitude}{direção} $045$: esse plano é provavelmente a falha, e o deslizamento sobre ele é \omterm{def:b1:faults-folds-fabrics:fault-types}{sinistral}.
\end{solution}

\begin{solution}{pb:b1:earthquake-source:1}
\textbf{1.} Os intervalos são $1290 - 1050 = 240$, $1480 - 1290 = 190$ e
$1712 - 1480 = 232$ anos; sua média é \qty{220.7}{yr}.

\textbf{2.} O deslocamento médio é $(7.0 + 6.0 + 5.5 + 7.5)/4 = \qty{6.5}{m}$, e
$T = 6.5/0.030 = \qty{216.7}{yr}$, próximo do intervalo médio.

\textbf{3.} Os deslocamentos dos três últimos terremotos, $6.0 + 5.5 + 7.5 =
\qty{19.0}{m}$, acumulados entre 1050 e 1712, em 662 anos: $\qty{19.0}{m}
/\qty{662}{yr} = \qty{28.7}{mm.yr^{-1}}$, \qty{4}{\percent} abaixo da geodésica \qty{30}{mm.yr^{-1}}. As duas concordam dentro da incerteza das datas e dos deslocamentos: o segmento libera quase toda a sua carga em seus grandes terremotos.

\textbf{4.} $2025 - 1712 = 313$ anos, e o déficit é
$\qty{0.030}{m.yr^{-1}}\times\qty{313}{yr} = \qty{9.39}{m}$.

\textbf{5.} O déficit excede cada deslocamento da trincheira, o maior sendo \qty{7.5}{m}, e o tempo desde 1712 excede o maior intervalo, 240 anos. Nesse sentido o segmento está atrasado; isso não significa que o terremoto seja iminente, apenas que a falha já retém mais deslizamento do que jamais liberou de uma só vez no registro.

\textbf{6.} Os intervalos vão de 190 a 240 anos e os deslocamentos de
\qtyrange{5.5}{7.5}{m}: o ciclo não é periódico, e a resistência da falha no próximo terremoto é desconhecida. A trincheira dá uma média e uma dispersão, a partir das quais uma probabilidade pode ser estimada, não uma data.

\textbf{7.} $A = \qty{250}{km}\times\qty{15}{km} = \qty{3750}{km^2} =
\qty{3.75e9}{m^2}$.

\textbf{8.} $M_0 = \num{3e10}\times\num{3.75e9}\times\num{9.39} =
\qty{1.06e21}{N.m}$.

\textbf{9.} $M_0 = \num{3e10}\times\num{3.75e9}\times\num{7.5} =
\qty{8.44e20}{N.m}$, logo $M_w = \tfrac23(20.926 - 9.1) = 7.88$.

\textbf{10.} A razão é $9.39/7.5 = 1.25$, que é $\tfrac23\log_{10}1.25 =
0.07$ de magnitude.

\textbf{11.} A metade da área reduz o momento à metade e baixa a magnitude em
$\tfrac23\log_{10}2 = 0.20$: $M_w = 7.75$.

\textbf{12.} $\mu A V = \num{3e10}\times\num{3.75e9}\times\num{0.030} =
\qty{3.38e18}{N.m}$ por ano.

\textbf{13.} $b = 1$. $N(5) = 10^{4.3 - 5} = 0.20$ por ano e $N(6) = 0.020$ por
ano.

\textbf{14.} $N(7) = 0.0020$ por ano: um a cada 500 anos; $N(8) = 0.00020$ por
ano: um a cada \num{5000} anos.

\textbf{15.} A lei dá um terremoto de magnitude 7.95 ou mais a cada
$1/10^{4.3 - 7.95} = \num{4500}$ anos, a trincheira um a cada 220 anos, vinte
vezes mais frequentemente. A lei é ajustada aos muitos terremotos pequenos de todas as falhas da região; extrapolada para os grandes terremotos de um só segmento, ela os subestima muito. A recorrência de grandes terremotos em uma falha é tomada de seu próprio registro e de sua carga, não dos pequenos.

\textbf{16.} $\num{5.5e17}/\num{3.38e18} = 0.16$: os terremotos moderados
liberam \qty{16}{\percent} do momento que o segmento acumula.

\textbf{17.} Pela \cref{prop:b1:earthquake-source:moment-steps},
$10^{1.5\times(7.95 - 6)} = 10^{2.92} \approx 840$ vezes mais.

\textbf{18.} Com $b = 1$, o momento dos terremotos é dominado pelo maior: o que quer que façam os terremotos moderados, liberam uma parcela menor, aqui
\qty{16}{\percent}, e os \qty{84}{\percent} restantes, o déficit de deslizamento, esperam pelo grande terremoto.

\textbf{19.} $v(-x) = -v(x)$ para cada par, logo o perfil é antissimétrico. Longe da falha as velocidades tendem a $\pm V/2 = \pm\qty{15}{mm.yr^{-1}}$.

\textbf{20.} $D = x/\tan(\pi v/V) = 15/\tan(\pi\times7.5/30) = 15/\tan\ang{45} =
\qty{15}{km}$: na \omterm{def:b1:earthquake-source:locking-depth}{profundidade de travamento} a velocidade é $V/4$.

\textbf{21.} Em \qty{5}{km}: $D = 5/\tan(\pi\times3.1/30) = 5/\tan\ang{18.6} =
\qty{14.9}{km}$; em \qty{40}{km}: $D = 40/\tan\ang{69.6} = \qty{14.9}{km}$. As velocidades são dadas com \qty{0.1}{mm.yr^{-1}}, e esse arredondamento sozinho desloca $D$ em cerca de \qty{0.1}{km}; com dados reais, o ruído de cada estação seria maior, e a \omterm{def:b1:earthquake-source:locking-depth}{profundidade de travamento} seria ajustada a todas as estações de uma vez.

\textbf{22.} A \omterm{def:b1:earthquake-source:locking-depth}{profundidade de travamento}, \qty{15}{km}, iguala a profundidade dos terremotos mais profundos. A parte travada da falha é sua parte frágil, a \omterm{def:b1:stress-strain-rheology:strength-envelope}{camada sismogênica}, onde o atrito a mantém travada entre os terremotos e onde ela falha neles; abaixo, a rocha mais quente rasteja (\cref{ch:b1:stress-strain-rheology}) e a falha desliza de modo constante, que é o deslizamento à taxa $V$ abaixo de $D$ do modelo.

\textbf{23.} $V/(\pi D) = \num{0.030}/(\pi\times\num{15000}) = \num{6.4e-7}$
por ano.

\textbf{24.} Com $M_0 = \qty{1.056e21}{N.m}$, $\log_{10}M_0 = 21.024$ e
$M_w = \tfrac23(21.024 - 9.1) = 7.949$. Com 0.01, o terremoto que liberaria todo o déficit do segmento tem a \omterm{def:b1:earthquake-source:mw}{magnitude de momento}
\textbf{7.95}.
\end{solution}
